\section{执行摘要}
\label{sec:summary}

\subsection*{一句话定义}

\keyword{Vulca} 是一款面向 AI agent 的 \keyword{agent-native} 图像处理 SDK ——
当所有团队都在卷 agent 的"大脑"（规划、推理、决策）时，Vulca 定位为 agent 的\keyword{"双手"}：
语义图层拆分、区域重绘、文化属性分析、合成导出，通过 \keyword{MCP（Model Context Protocol）}
把像素级执行能力暴露给 Claude Code、Cursor 等 agent 宿主，
让 agent 从"描述创作意图"升级到"真正完成可交付的视觉产出"。

\subsection*{当前进度（截至 2026-04-23）}

\begin{itemize}
  \item \fact{产品版本}{v0.17.5（PyPI 已发布，Apache 2.0 开源）}
  \item \fact{技术基线}{21 个 MCP 工具 \textperiodcentered{} 高覆盖率测试套件（CI 绿灯率 $\ge$ 99\%，详见 GitHub repo CI badge）\textperiodcentered{} Phase 0 和 Phase 1 均已完成}
  \item \fact{已上线 skill}{\code{/decompose}（语义分层）\textperiodcentered{} \code{/visual-brainstorm}（视觉意图 brief）\textperiodcentered{} \code{/visual-spec}（可评审设计文档）}
  \item \fact{本地化链路}{Apple Silicon 上 ComfyUI + Ollama + EVF-SAM 端到端跑通，无需任何云端密钥}
  \item \fact{模型生态}{已完成 Google Gemini 全套模型集成（2.5 VLM 视觉理解 + 3.1 image preview 生成系列）与 OpenAI、本地 ComfyUI 的多 provider 适配；Gemini 与 Anthropic 研究/创业 API credits 已到位，短中期算力成本可控}
  \item \fact{文化评估}{覆盖中国工笔、中国写意、日本传统、伊斯兰几何、西方学院、南亚、非洲、摄影、水彩、品牌设计、UI/UX、当代艺术共 12 个具体艺术/设计传统 + 1 个通用基线（合计 13 种）的 L1–L5 可引用评价准则}
  \item \fact{学术背书}{承接 ACM MM 2026 EmoArt Challenge（2026-06-10 提交），Track 1 基于 Vulca 技术栈}
  \item \fact{早期社区信号}{GitHub 插件生态（\code{vulca-org/vulca-visual-agent-plugin}）\textperiodcentered{} dev.to + 小红书双语种子内容已发布 \textperiodcentered{} 流量与下载数据周期性复盘}
\end{itemize}

\subsection*{市场机会一句话}

\keyword{全球 agent 浪潮 $\times$ 多模态 SDK 的"编辑 / 分层 / 评估"空白 $\times$ 中国文化资产数字化与自主可控政策窗口} ——
三股趋势交汇处，Vulca 是唯一做\keyword{"agent-native 像素执行"}的开源栈，
同时具备\keyword{国际生态嵌入}能力（Claude Code 插件）与\keyword{国内自主可控}属性（本地链路、文化评估、Apache 2.0）。

\subsection*{商业模式（Open-Core 三层漏斗）}

\begin{enumerate}
  \item \keyword{开源社区导流}：Apache 2.0 SDK + MCP 工具集免费，吸引 agent 开发者与独立设计师，建立品牌与生态粘性；
  \item \keyword{托管服务与文化评估订阅}：为团队客户提供合规 SaaS（SSO、私有数据、SLA）和独立的文化属性分析 API（按次/月订阅）；
  \item \keyword{行业本地化部署}：面向博物馆、出版社、品牌设计机构、高校艺术研究中心，提供\keyword{完全离线}的本地部署 + 年度许可 + 定制实施。
\end{enumerate}

\subsection*{本轮融资诉求}

\begin{itemize}
  \item \fact{阶段}{天使轮（Angel Round）}
  \item \fact{募资金额}{¥ 600 万（区间 ¥ 500–800 万；按 18 个月跑道、7 人团队反推）}
  \item \fact{估值口径}{post-money ¥ 3500 万（对标同阶段 agent 与文化科技方向天使轮估值中位数）}
  \item \fact{资金用途}{约 80\% 用于团队扩展（工程、算法、BD），约 20\% 用于行业客户落地与市场教育；12–18 月跑道}
  \item \fact{已覆盖资源}{Gemini 全套 + Anthropic 研究级 API credits 已到位，短期算力几乎零现金成本，资金全部用于人与市场}
\end{itemize}

\subsection*{为什么是现在投/现在报}

\begin{itemize}
  \item \keyword{技术已验证}：21 个 MCP 工具 + 高覆盖率测试套件 + 3 个端到端 skill 全部 shipped，不是 PPT 项目；
  \item \keyword{生态位稀缺}：agent-native + 像素执行层，头部厂商（OpenAI、Adobe、Runway）选择了封闭"大脑"路线，HuggingFace 生态只给零件不给工作流，此中间层无直接竞争者；
  \item \keyword{双赛道受益}：既吃国际 agent 浪潮（Claude Code / Cursor 插件生态）红利，又契合国内文化资产数字化 + 自主可控政策窗口；
  \item \keyword{学术+产业双轮}：ACM MM 2026 EmoArt Challenge 提交在即，学术成果直接转化为产品能力，构成持续的\keyword{研究含量}证明。
\end{itemize}
